\documentclass[usenatbib]{mnras}

\usepackage{amsmath}
\usepackage{graphicx}

\title[SCM-BH: Regime transition in AGN jet collimation]{SCM-BH: Empirical regime transition in AGN jet collimation from the MOJAVE catalogue}

\author[S. C\'amara Madrid]{Sergio C\'amara Madrid}

\begin{document}

\maketitle

\begin{abstract}
We present an empirical analysis of AGN jet opening angles using the MOJAVE catalogue
(VizieR J/MNRAS/468/4992). Applying a threshold at apparent brightness-temperature
ratio $r_{15} \approx 10$, we identify a statistically robust regime separation
between a LOW ($r_{15} < 10$, $N = 274$) and a HIGH ($r_{15} \geq 10$, $N = 86$)
sub-sample. A Kolmogorov--Smirnov test yields $p \approx 5.6 \times 10^{-16}$,
and the Spearman rank correlation within the HIGH regime is
$\rho \approx -0.331$ ($p \approx 0.0019$), indicating progressive collimation.
No underlying physical mechanism is assumed.
\end{abstract}

\begin{keywords}
galaxies: jets -- quasars: jets -- BL Lacertae objects: general
\end{keywords}

\section{Introduction}
\label{sec:intro}

% Placeholder -- manuscript in preparation.

\section{Data}
\label{sec:data}

The source catalogue is the MOJAVE survey, available via VizieR (J/MNRAS/468/4992).
The analysis uses columns \texttt{r15} (ratio of brightness temperatures) and
\texttt{alphaApp15} (apparent jet opening angle in degrees).

\section{Analysis}
\label{sec:analysis}

\subsection{Regime separation}

A threshold at $r_{15} = 10$ ($\log r_{15} = 1$) divides the sample into two regimes.

\subsection{Statistical tests}

\begin{itemize}
  \item KS two-sample test: $p \approx 5.6 \times 10^{-16}$.
  \item Spearman rank correlation (HIGH regime): $\rho \approx -0.331$, $p \approx 0.0019$.
\end{itemize}

\section{Results}
\label{sec:results}

% See figures fig1\_transition.png and fig2\_high.png in results/.

\section{Discussion}
\label{sec:discussion}

% Manuscript in preparation.

\section{Conclusions}
\label{sec:conclusions}

We report a statistically significant regime separation in AGN jet opening angles
at $r_{15} \approx 10$, with progressive collimation in the HIGH regime.
The threshold was identified empirically and tested for stability.

\bibliographystyle{mnras}
\bibliography{references}

\end{document}
